기대값 기반의 결정 이론은 선택의 기대값이 0 이상이라면 그 선택을 하는 것이 합리적이라고 판단한다. 본 논문에서 주장하려는 종류의 합리성은 인간은 본전을 되찾을 확률이 50\% 이상인 게임을 하는 것이 합리적이라고 판단한다는 것이다.

일반적인 확률 아래의 게임의 경우, 기대값 기반의 합리성과 본전을 찾을 확률을 고려한 합리성이 일치한다. 그러나 희박한 확률 아래의 게임에서는 기대값 기반의 합리성만 이용한다면 상식에 어긋나는 결과를 도출할 수 있다.

희박한 확률 아래의 게임에서는 게임의 기대값이 0 이상이더라도 efficient time 이내에 참가자가 본전을 찾을 확률이 낮다는 것이 그 이유이다.

희박한 확률 아래의 게임을 참가하는 인간은 기대값 뿐만 아니라 본전을 찾을 확률 또한 고려한다고 볼 수 있다.
